% =============================================================================
% 1.1 研究背景与意义
% 草稿版本：v0.1 (2026-05-05)
% 字数：约 1500 字
% 风格规则：v0.2 风格（无 \textbf / \textit / 引例括号 / 双引号 / 破折号）
% 引用：\cite{ref2, ref5, ref10, ref18, ref25}
% 依赖宏包：booktabs / tabularx（主稿已加载）
% =============================================================================

\subsection{研究背景与意义}
\label{sec:background}

气象预报数据是支撑国民经济与社会发展的重要基础信息资源。随着全球
气候变化加剧，极端天气事件的发生频率与强度持续攀升，精准获取与
深度分析气象预报信息已成为防灾减灾、农业生产、能源调度、交通运输
等众多领域的迫切需求 \cite{ref2, ref5}。据中国气象局公布的统计
数据，我国每年因气象灾害造成的直接经济损失占国内生产总值的 1\%
至 3\%，气象服务的精准化与智能化对降低灾害损失具有现实意义。

然而，当前气象预报数据的获取与分析面临两类深层次技术鸿沟，一类
来自数据接入侧的接口异构性，另一类来自数据解读侧的认知差距，二
者在不同环节同时阻碍着气象数据的智能化利用。

数据接入侧的异构性鸿沟体现为多源数据规范的不统一。气象数据来源
于中国气象局、欧洲中期天气预报中心、美国国家海洋和大气管理局
以及和风天气、Open-Meteo 等多种公开商业平台，不同来源的观测
数据、数值预报产品、卫星遥感资料在数据格式、时空分辨率、要素
类型上差异巨大 \cite{ref25}。专业用户需同时掌握多种 API 接口
规范、字段命名约定与错误码体系才能完成跨平台的数据检索。例如
和风天气以 \texttt{LocationID} 作为地点标识符并以 \texttt{code}
字段返回状态码，而 Open-Meteo 直接以经度纬度作为参数并以
\texttt{error} 字段返回错误描述。这种接口层异构性对终端用户构成
了显著的工程门槛。

数据解读侧的语义鸿沟体现为数值表达与决策语义之间的形态错位。气象
数据以高度结构化的数值形式存在，例如 24 小时累积降水量为
\SI{52.3}{mm}、10 米高度风速为 $\SI{14.1}{m/s}$，而决策语义却以
自然语言形式表达，例如暴雨预警、不宜户外作业、建议携带雨具。
普通用户难以从大量原始数值中提炼出具有决策参考价值的判断。这一从
数值数据到决策语义的认知鸿沟使得气象预报数据的潜在价值未能被
充分释放。

与此同时，以大语言模型为核心的 AI 智能体技术正在迅速发展。Xi 等
\cite{ref10} 在综述中系统梳理了基于大语言模型的智能体在感知、推理、
规划、工具使用等方面的能力框架。Yao 等 \cite{ref18} 提出的 ReAct
范式让大语言模型能够在推理过程中与外部工具交互执行动作，为将大
语言模型应用于需要外部数据获取与计算的实际任务提供了可行路径。
李特等 \cite{ref2} 与代刊等 \cite{ref5} 分别探讨了 AI 智能体与
大语言模型在天气预报业务中的应用前景，指出利用大语言模型的自然
语言理解与推理能力来桥接用户需求与气象数据之间的鸿沟，是气象
信息服务智能化的重要方向。

基于上述背景，本文针对气象数据接入侧的异构性问题与数据解读侧的
语义鸿沟问题，设计并实现一个基于 AI 智能体的气象预报数据检索与
分析系统。系统以大语言模型为推理引擎，采用 ReAct 范式构建单
智能体架构，通过工具学习自动调用气象 API 完成数据检索，并融合
代码生成与语义桥接机制实现确定性的统计分析与自适应的报告生成。
本文研究意义可从工程实践、方法论沉淀与评测体系三个层面归纳。

在工程实践层面，本文为面向终端用户的气象信息服务智能化提供了一
套从自然语言查询到决策报告的完整技术链路，把原本需要专业人员
逐步操作的多源 API 调用、数据筛选、统计计算、术语解读、规范引用
等繁琐工作一次性自动化，把气象数据获取的工程门槛从 API 编程
水平降低到自然语言对话水平。

在方法论沉淀层面，本文提出的双通路 RAG 架构与代码生成加受限沙箱
方案在气象领域系统化地解决了大语言模型的数值与引用幻觉问题。
该方法论的核心思路是把高风险任务从大语言模型内部转移到外部确定
性组件，让大语言模型只承担擅长的语言组织职责，避免在数值边界、
权威引用、统计计算等模型不擅长的环节强行使用其内部知识。这一
思路在其他垂直领域的智能体设计中具有可迁移性。

在评测体系层面，本文构建了 5 套覆盖意图识别、检索、桥接、代码
执行、端到端共 214 条用例的多层评测体系，并在系统开发过程中
完成了 3 次知识库的人工核对与 24 处错误修正。该评测过程为本文
所设计方法在意图识别、语义桥接等关键环节的工程可靠性提供了
可复现的量化记录，所产生的评测集与核对表也可作为后续相关研究的
对照基线。

本章后续内容首先综述与本课题密切相关的国内外研究进展，进而阐述
本文的研究内容、研究方法与论文特色，最后给出全文的组织结构。

